\begin{textsample}{POS Dim 3 – pb – Score 20.00 – pb/pb\_inf\_14.txt}  \label{ex:f3_pos_112}
CAMPO DE EXPERIÊNCIA : CORPO, \textbf{GESTOS} E MOVIMENTOS OBJETIVOS DE APRENDIZAGEM E DESENVOLVIMENTO Movimentar as partes do corpo para exprimir corporalmente emoções, necessidades e \textbf{desejos}.
Experimentar as possibilidades corporais nas \textbf{brincadeiras} e interações em ambientes acolhedores e desafiantes. \textbf{Imitar} \textbf{gestos} e movimentos de outras \textbf{crianças}, \textbf{adultos} e animais.
Participar do \textbf{cuidado} do seu corpo e da promoção do seu bem-estar.
Utilizar os movimentos de preensão, \textbf{encaixe} e lançamento, ampliando suas possibilidades de \textbf{manuseio} de diferentes materiais e objetos. ( BNCC, 2017 ) SUGESTÕES DE VIVÊNCIAS QUE FAVORECEM A CONSTRUÇÃO DE EXPERIÊNCIAS PELOS \textbf{BEBÊS} NESTE CAMPO As práticas realizadas na Instituição de Educação \textbf{Infantil} devem possibilitar : - Exploração dos espaços por meio de movimentos como engatinhar, arrastar -se, andar, descer, subir, transpor obstáculos, bem como outros ; - Expressão por meio de \textbf{gestos}, movimentos corporais, \textbf{balbucios} e da fala para se comunicar com seus pares e \textbf{adultos} ; - Socialização de momentos de \textbf{brincadeiras} com outras \textbf{crianças} e \textbf{adultos} ; - Exploração, conhecimento e reconhecimento de diferentes partes do corpo ; - \textbf{Imitação} de \textbf{gestos} de outras pessoas e reprodução de diferentes ruídos e \textbf{sons} ; - \textbf{Acompanhamento} do ritmo das músicas com movimentos corporais ; - \textbf{Imitação} de animais, personagens e estados de ânimo ; - Manipulação de objetos de diferentes formas, \textbf{pesos}, \textbf{texturas}, tamanhos, que produzam \textbf{sons} etc., usando movimentos como pegar, largar, levar à \textbf{boca}, chutar, \textbf{empilhar}, lançar em várias direções e de diferentes modos ; - Realização de atividades expressivas que busquem a criação de movimentos coordenados ; - Momentos de interações com outros grupos, com diferentes possibilidades de movimento em ambientes diversos ; - Uso de espelhos para se reconhecer e aos seus \textbf{colegas}, visualizar seus \textbf{gestos} e movimentos ; - Valorização de diversas expressões culturais próprias ao universo dos \textbf{bebês}, especialmente aquelas que inspirem movimentos do seu corpo.

% matched lemmas: acompanhamento, adulto, balbucio, bebê, boca, brincadeira, colega, criança, cuidado, desejo, empilhar, encaixe, gesto, imitar, imitação, infantil, manuseio, peso, som, textura
\end{textsample}
